% ======================================================================
% Section — The main theorem, stated conditionally
% ======================================================================
% Placed after the introduction. The gloss keeps the statement
% self-contained; the derivation and calibration are in sec "scale".

\section{The main theorem, stated conditionally}\label{sec:SC}

The closures of \S\ref{sec:rigid}--\S\ref{sec:higherrung} settle the
covering obstruction at the base rung, at the composite micro-cases, and
at the critical cells of the higher rungs. What remains is the frontier:
the $(1K,5U)$ cell at $T=8$ and its generalizations $(1K,T{-}3U)$ at
every $T\ge 8$, the cells whose five-and-more-difference censuses are too
large to enumerate and whose closure must therefore be structural. This
section states the program's central result in what we judge its only
honest form: a single named conditional theorem, preceded by a ledger of
hypotheses in which every input carries its status---proved, measured,
or derived from a disclosed fitted constant. The theorem is the
consolidation of the conditional statements scattered through the
cascade analysis of \S\ref{sec:scale}; stating it once, with named
hypotheses, is the point of this section.

The statement needs the higher-rung vocabulary of
\S\ref{sec:higherrung}, which follows later; we give the gloss here and
defer all derivations. At $N=p^2$ with $p=Tk+\epsilon$ prime, a covering
family consists of \emph{kernel} members (speeds divisible by $p$) and
\emph{unit} members (speeds coprime to $p$); a cell $(aK,bU)$ collects
the families with $a$ kernel and $b$ unit members. The kernel covers the
fibers of a kernel bad set and leaves $|U|=(T-2)k+\epsilon-1$ uncovered
fibers $F_j\cong\Z_p$; on each, the $b$ unit footprints are arithmetic
progressions which must cover $\Z_p$ (Lemma~\ref{lem:rungT}). The
\emph{census normal form} pins the normalizer arc's difference to $1$
and takes canonical differences in $[2,D]$, $D=(p-1)/2$; for the
frontier cell, $m=T-3$ arcs carry a two-value size window
$\{2k,2k+1\}$ or $\{2k+1,2k+2\}$ and an \emph{overlap budget}
$\mathrm{ov}=\Sigma s-p\approx(T-6)k+O(1)$. A family is
\emph{census-admissible} if it admits at least one covering placement in
the normal form; the \emph{fiber cascade} is the iterated filter of
\S\ref{sec:scale-model}, which retains the anchor placements surviving
the uncovered fibers one by one and kills the family when the survivor
set empties.

\begin{theorem}[scaling closure; conditional]\label{thm:SC}
Assume the hypotheses $\mathrm{NF}$, $\mathrm{Hq}$, $\mathrm{H1g}$,
$\mathrm{H2r}$ of Table~\ref{tab:ledger} (statuses as in the table;
$\mathrm{Hq}$ subsumes the former pairwise input $\mathrm{H2m}$ as its
one-step case), together with the committed equivalence $\mathrm{EQ}$
of the same table. Let $T\ge 8$, let $p=Tk+\epsilon$ be prime with
$k\ge 2$ and $0\le\epsilon<T$ (so $k=\floor{p/T}$ is canonical; $p\ge
p_0(T)$, the least such prime, $p_0(8)=17$), and let the frontier cell
$(1K,T{-}3U)$ at $N=p^2$ be given, with $m=T-3$ unit arcs, the
two-value size window of Lemma~\ref{lem:rungT}, and overlap budget
$\mathrm{ov}=\Sigma s-p=(T-6)k+O(1)$. Then:
\smallskip
\emph{(i)} every census-admissible family in the cell dies in the fiber
cascade within
\begin{multline*}
C(T,k)\;=\;
\frac{\log\!\bigl(c_1\,\mathrm{ov}^{\,m-1}\bigr)}
     {(m-1)\,\log\!\bigl(p/\mathrm{ov}\bigr)-\log(c_1\Gamma)}
+\;\frac{2\Lambda_0\,(T-2)}{T}\;+\;1\;+\;O(1/k)
\end{multline*}
fibers, where $\Gamma$ is the uniform constant of $\mathrm{Hq}$; the
denominator is uniformly positive under $\mathrm{H1g}$'s margin clause,
and its $k\to\infty$ form $(m-1)\log(T/(T-6))-\log(c_1\Gamma)$ gives
the $T$-flat plateau $\approx 9$, $9.2$, $10.1$ at $T=8,9,10$ (at
$k=2,3,3$)---with the never-exceeded worst case
\[
(T-4)+N_{\mathrm{res}}+O(1)\;\le\;T+3,
\qquad
N_{\mathrm{res}}=1+\tfrac{2\Lambda_0(T-2)}{T};
\]
\smallskip
\emph{(ii)} since $C(T,k)\ll|U|=(T-2)k$ for every $T\ge 8$, $k\ge 2$,
with a surplus growing linearly in $k$ (conditional on the ledger: the
depth bound is what (i) provides; the surplus language refers to the
fiber budget, not to an independently proved decay rate), the cascade
exhausts the fiber budget: no census-admissible family admits a
covering placement valid on all of $U$, and hence
\smallskip
\emph{(iii)} the cell $(1K,T{-}3U)$ at $p^2$ is empty; and by the
lossless reformulation equivalence, $\mathrm{LRC}$ holds at that rung.
The measured range of the calibration is $T=8,9,10$ at $p\le 73$, with
committed ground-truth cells at $p\le 61$ for $T\in\{7,8\}$; the
statement at primes beyond the calibration is exactly what the
hypotheses purchase.
\end{theorem}

\begin{proposition}[the Hq-induction; the derivation step]\label{prop:Hq}
Let the frontier cell and notation be as in Theorem~\ref{thm:SC}, and
assume $\mathrm{Hq}$ with constant $\Gamma$ and the volume law
$\mathrm{H1g}$ with $|\mathrm{Sol}(f,\mathrm{sz})|\le
c_1\,\mathrm{ov}^{\,m-1}$ for every admissible family and size tuple.
Write $S_i$ for the survivor set after $i$ processed fibers,
$A_i=\lambda_i^{-1}(\mathrm{Sol}_i-\tau_i)$ for the $i$-th fiber's
dilate constraint, $d=m-1$, $B=c_1\,\mathrm{ov}^{\,d}$, and
$r=\Gamma B/p^{\,d}$. Then $|S_0|\le B$, each step contracts,
$|S_i|\le r\,|S_{i-1}|$, and consequently
\[
|S_q|\;\le\;B\,r^{\,q},
\qquad\text{so the cascade is empty whenever}\qquad
q\;>\;\frac{\log B}{(m-1)\,\log(p/\mathrm{ov})-\log(c_1\Gamma)}.
\]
The denominator is uniformly positive iff
$c_1\Gamma<\bigl(p/\mathrm{ov}\bigr)^{m-1}$ with a margin uniform in
$(T,k)$; in the $k\to\infty$ limit $p/\mathrm{ov}\to T/(T-6)$, and the
uniform condition is the threshold clause of $\mathrm{H1g}$.
\end{proposition}

\begin{proof}
$S_0$ is the family's initial placement set, bounded by the volume law.
For each $i$, $S_i=S_{i-1}\cap A_i$ with $A_i$ the $i$-th fiber's
constraint set, so $\mathrm{Hq}$ applied to the admissible history
$S_{i-1}$ gives $|S_i|\le\Gamma\,|S_{i-1}|\,|\mathrm{Sol}_i|/p^{d}\le
r\,|S_{i-1}|$ using $|\mathrm{Sol}_i|\le B$. Induction gives
$|S_q|\le B\,r^{q}$, and $|S_q|<1$ forces emptiness. The resonance tail
(the weak-cut fibers with $\bar\lambda\le\Lambda_0$, each of which may
fail to contract) adds $N_{\mathrm{res}}=1+2\Lambda_0(T-2)/T$ fibers
before the contraction count begins, which is the tail of $C(T,k)$.
\end{proof}

\paragraph{What changed in this revision, and why.}
The theorem's predecessor assumed an unquantified ``sequential
composition'' input (the former $\mathrm{H0}$ row: $C_i\approx
\bar\gamma^{\,i}|\mathrm{Sol}|^{i+1}/p^{(m-1)i}$) and a \emph{pairwise}
dilate-sparsity input (the former $\mathrm{H2m}$). As an audit of this
paper observed, a recurrence for the \emph{independent-extension model}
and a bound on \emph{unconditional} pair intersections do not compose
into a bound on the \emph{iterated, history-dependent} intersections the
actual cascade performs: pairwise sparsity does not control
higher-order correlations, and three sets can have acceptable pairwise
overlaps while sharing a structured common core. The repair is minimal
and explicit: the multi-fiber survivor inequality $\mathrm{Hq}$ (above)
replaces both, with $\mathrm{H2m}$ retained as its one-step case and
$\mathrm{H0}$'s measurements retained as its calibration support. The
derivation from the ledger is then exactly
Proposition~\ref{prop:Hq}'s induction---complete, and conditional on
named open inputs only.
\emph{The third revision (v3) changes no theorem, lemma, or hypothesis
of this monograph}: it updates the zonotope-strand record of
\S\ref{sec:scale} for the companion paper's closure of the last open
certification row (the $n=5$, $m=10$ deepest-hole value, now exactly
$7/22$ by a Lipschitz box certificate), harmonizes the harmonic
lower-bound record with the companion's audited $4183/12288$ witness,
and records the version marker. \emph{The fourth revision (v4) again
changes no theorem, lemma, or hypothesis}: it updates the
zonotope-strand record for the companion's third certification
session---the harmonic instance's deepest hole closed \emph{exactly}
($\rho(1,2,3,4,5)=16/47$, by a balance-family witness whose three tied
pair-dips obey a weighted identity, certified by branch and bound with
two new box certificates), and the two $n=6$ boundary rows resolved as
strict refutations ($\rho\ge17/47>5/14$ at $m=12$,
$\rho\ge455/1269>5/14$ at $m=9$)---and records the version marker.

\paragraph{What is conditional and what is not.}
The reduction chain \emph{(i)}$\to$\emph{(ii)}$\to$\emph{(iii)} is
proved from committed machinery: the two-fiber reduction
(Lemma~\ref{lem:twofiber}: the cascade step \emph{is} a
dilate-intersection of census solution sets), the drift law
(Lemma~\ref{lem:rungT}, machine-verified $3{,}528/3{,}528$), the
$T$-free fiber machinery, and the lossless equivalence of
Corollary~\ref{cor:tau}. The arithmetical content---that the per-fiber
cut beats the pool within $C(T,k)$ fibers---is proved \emph{given} the
ledger's inputs, by the induction of Proposition~\ref{prop:Hq} (in
this revision an explicit theorem-level step; its predecessor leaned
on the independent-extension recurrence without the quantified Hq). None of those inputs is silently established, and none
is silently weakened: that is the content of Table~\ref{tab:ledger}. The
$pq$ side of the composite frontier is carried by the row translation
(Theorem~\ref{thm:R}) with the same input structure; Theorem~\ref{thm:SC}
is stated for the $p^2$ side only.

\paragraph{The hypothesis ledger.}
Each row of Table~\ref{tab:ledger} names one input, its precise form,
its status, and its evidence or attack route. The rows are
\emph{conditional inputs}, not evidence toward the theorem's validity:
three are open hypotheses with measured support ($\mathrm{Hq}$,
$\mathrm{H1g}$, $\mathrm{H2r}$, the first sharpened and quantified in
this revision); one is the measured one-step base case
($\mathrm{H2m}$); the others are proved, machine-validated structure
($\mathrm{NF}$) and a committed theorem ($\mathrm{EQ}$). The statuses
are stated once here and are not restated with each citation.

{\small
\begin{longtable}{@{}p{0.075\textwidth}p{0.355\textwidth}
                   p{0.155\textwidth}p{0.315\textwidth}@{}}
\caption{The hypothesis ledger of Theorem~\ref{thm:SC}. Statuses:
\emph{OPEN} (a conditional input with measured support), \emph{proved}
(with the verification record), \emph{measured} (base case with tested
range), or \emph{committed}. The substantive open inputs are Hq, H1g,
and H2r; every row is an input the theorem assumes, not evidence that
the theorem holds.}\label{tab:ledger}\\
\toprule
& hypothesis (precise form) & status & evidence / route \\
\midrule
\endfirsthead
\multicolumn{4}{@{}l@{}}{\itshape Table~\ref{tab:ledger}, continued}\\
\toprule
& hypothesis (precise form) & status & evidence / route \\
\midrule
\endhead
\bottomrule
\endfoot
$\mathrm{NF}$ &
Census normal-form completeness: every covering family at an uncovered
fiber is census-admissible (the normal form---normalizer difference
$1$, canonical differences, two-tier sizes---loses no coverings). &
proved in structure, machine-validated &
The footprint and arc derivations are $T$-free and verified
$957/957$ against direct $\Z_{p^2}$ enumeration; the $T=7$
census-restricted search equals brute force at the boundary prime
$p=5$ (the $8$-covering control); the $T=8$ zoo census is the measured
object. \\
$\mathrm{Hq}$ &
Multi-fiber survivor inequality (this revision; replaces the former
unquantified H0 and subsumes H2m as its one-step case): for every
admissible cascade history $S_{i-1}$ and next fiber $i$,
$|S_{i-1}\cap A_i|\le\Gamma\,|S_{i-1}|\,|\mathrm{Sol}_i|/p^{m-1}$,
where $A_i=\lambda_i^{-1}(\mathrm{Sol}_i-\tau_i)$, with $\Gamma$
uniform over families, size tuples, dilations, translations, and prior
surviving fibers. &
OPEN (measured benign in range) &
Depths $\le 4$ at the committed cells ($764$-pair record); $\le 11$
across the three calibration rungs; exit time $\le 2$ exhaustive over
admissible families at $T=7$ and $T=8$-critical; the $i=1$ case is the
measured H2m. Route: tangency codimension $\ge 1$; the pairwise
measurements are the base case, the iterated conditional bound is the
open content. \\
$\mathrm{H1g}$ &
Distinct-family volume law with uniform margin (sharpened in this
revision): $|\mathrm{Sol}(f,\mathrm{sz})|\le
c_1\,\mathrm{ov}^{\,m-1}$ for every $\pm$-distinct family $f$ and
admissible size tuple, with $c_1\le\bigl(T/(T-6)\bigr)^{\,m-1}
e^{-\eta}$ for some $\eta>0$ uniform in $(T,k,f,\mathrm{sz})$---i.e.,
the constant beats the threshold
$c_1=\bigl(T/(T-6)\bigr)^{\,m-1}$ ($256/243/244.1$ at $T=8/9/10$,
where the depth denominator of Proposition~\ref{prop:Hq} vanishes) by
a fixed multiplicative margin. The earlier phrasing ``$c_1=O(1)$
uniform, equivalently beats the trivial bound by a factor growing in
$k$'' is retained as calibration language, not as an equivalent
hypothesis. &
OPEN (measured, margin growing in $k$; chain class proved; three
attack routes closed, the automaton route by theorem) &
At $T=8$: worst ratio $|\mathrm{Sol}|/\mathrm{ov}^4$ of
$0.214/0.083/0.013$ at $p=17/19/29$ (chain-class benchmark
$2.17/3.61/2.03$); measured margins over the trivial box
$95\times/1{,}207\times/3{,}803\times$. Remaining distance: the strand
census (\S\ref{sec:scale-h1})---necessary-condition and
marginal-rigidity routes measured closed, the coordinates fixed by
Lemma~\ref{lem:char10}, and the transfer-matrix/automaton route closed
by Lemma~\ref{lem:monomial}: the instrument computes the census exactly
(its cyclic trace \emph{is} the survivor count) and consumes H1g rather
than proving it, every spectral quantity in the framework being exactly
$0$ or $1$. \\
$\mathrm{H2m}$ &
Mid-region dilate-sparsity---now the $i=1$ (one-step, unconditional)
case of $\mathrm{Hq}$, retained for its stronger measured record:
$|\mathrm{Sol}_0\cap\lambda^{-1}(\mathrm{Sol}_j-\tau)|\le
\gamma_0\,|\mathrm{Sol}_0|\,|\mathrm{Sol}_j|/p^{m-1}$ for every fiber
pair with $\bar\lambda>\Lambda_0$, $\gamma_0=O(1)$ uniform. &
OPEN (measured; proved for the chain class with loose constants) &
$764$ pairs at $5$ committed cells: mid-zone $\gamma$ median $0$
(instant kills), max $40.6$. Chain class proved
(Theorem~\ref{thm:chaindilate}) with constants $(m!)^2\cdot 4m\cdot
s_{\max}$---loose by orders of magnitude. Route: the
Lemma~\ref{lem:E} three-distance pair-line count. \\
$\mathrm{H2r}$ &
Resonance budget, now stated as an open assumption rather than
``derived'': the weak-cut fibers---those whose dilation ratio
$\lambda_i$ has small rational height, $\bar\lambda\le\Lambda_0$ in
the height function of \S\ref{sec:scale-model}---number $\le
2\Lambda_0(T-2)/T+1$, $k$-free; the $\lambda=-1$ fiber is never a cut
(proved, Lemma~\ref{lem:reflect}). The count's formula shape is
derived from the reflection lemma and a Gauss-circle-type heuristic;
$\Lambda_0\approx 4$ is data-chosen, not proved. &
OPEN (measured support; formula shape derived; the $\lambda=-1$ clause
proved) &
Lemma~\ref{lem:reflect} (the $x\mapsto-x$ involution; survival
fraction exactly $1$ at every committed family and cell); the count is
a Gauss-circle-type small-rational dilation count;
\emph{$\Lambda_0\approx 4$ is data-chosen}. \\
$\mathrm{EQ}$ &
Lossless reformulation: the cell emptiness at the rung implies
$\mathrm{LRC}$ there. &
committed theorem &
The pair-collapse equality
$M(V)=\max_{\{u,w\}}\tau(u,w)$ (Theorem~\ref{thm:equality},
machine-verified with a $2{,}870/2{,}870$ independent re-verification);
the fibration threshold-compatibility. \\
\end{longtable}
}

\paragraph{The free-parameter disclosure.}
Standing practice for this paper, applied to every constant entering
Theorem~\ref{thm:SC}: a claim is called \emph{derived} only when it is
derived outright, and a formula containing a fitted constant is labeled
as such. The constants are:
\begin{itemize}[leftmargin=2.2em,itemsep=2pt,topsep=2pt]
\item \emph{$c_1$ (H1g):} bounded by data only---$\le 0.214$ at $T=8$
  over the three census primes $p=17,19,29$, tightening with $k$
  (Table~\ref{tab:volume}); \emph{not proved}. The chain class, the only
  proved case, carries $c_1\approx 2.0$--$3.6$. The threshold the
  constant must beat is structural, not fitted: the trivial bound
  realizes $c_1=\bigl(T/(T-6)\bigr)^{\,m-1}$
  (Proposition~\ref{prop:threshold}), so the content of H1g is the
  margin over $\approx250$---which every measured constant carries and
  no proved constant yet does.
\item \emph{$\Lambda_0\approx 4$ (H2r and the resonance tail of $C$):}
  data-chosen---the $\bar\lambda$ zone in which the measured per-fiber
  cuts at the committed cells first exceed the random model by
  $\sim 10\times$. The formula
  $N_{\mathrm{res}}=1+2\Lambda_0(T-2)/T$ is therefore \emph{a derived
  formula containing one fitted constant}: partially derived, not
  derived outright.
\item \emph{$\gamma_0$ (H2m):} bounded by data ($\le 41$ in the
  non-resonant zone over the committed cells); proved only for the chain
  class, with loose constants.
\item \emph{$R$ (a junction-volume scale fitted at $4.2$ in an earlier
  formulation):} \emph{retired}. For the chain class it is replaced by
  the derived overlap budget $T-6$; for the thin families it is
  superseded by $c_1$ under H1g. No fitted $R$ remains in
  Theorem~\ref{thm:SC}.
\end{itemize}
The depth formula's remaining constants are structural: the leading
term is built from the overlap budget $T-6$ and the modulus $T$, and
the tail from Lemma~\ref{lem:reflect} (proved) plus the $\Lambda_0$
count above.

\paragraph{The sampled numbers, with intervals.}
Theorem~\ref{thm:SC}'s calibration rests on exhaustive censuses where
the cells are enumerable and on disclosed samples where they are not.
The $T=8$ saturation figures of Table~\ref{tab:pool}
($1.000/0.571/0.357/0.133/0.015$ at $p=17/19/23/29/31$) are exhaustive
censuses and carry no sampling error. The $T=9$ figures are
stride-sampled with full decisions per sampled tuple: at $p=29$,
saturation $0.903$, $95\%$ Wilson interval $[0.889,0.915]$
($n=2{,}005$), projected pool $1.39\times10^5$ with interval
$[1.37,1.41]\times10^5$; at $p=31$, saturation $0.490$,
$[0.468,0.511]$ ($n=2{,}002$), projected pool $1.18\times10^5$,
$[1.12,1.23]\times10^5$. Both pools are \emph{projected}---the
falling-factorial bound times the sampled fraction, at one prime
each---and the rung-monotone reading rests on two primes with
$\epsilon\in\{2,4\}$, whose modulation ($0.90$ against $0.49$) is large
and unmodeled. The $T=10$ saturation (extrapolated $\ge 0.9$ at the
$40\%$ budget from the rung-monotone pattern---two primes, unmodeled
$\epsilon$-modulation, no error control) is \emph{not measured}. Cascade-depth figures at the
committed cells are exact runs (median $2$, maximum $4$ over the zoo
random-$40$); the maxima $9/11/10$ are the measured plateau across
$k\le 9/8/3$ at the three rungs, quoted as measurements.

\paragraph{Scope levels, fixed once.}
Three different closures circulate in this paper and must not be
conflated: \emph{cells} (the $T=7$ and $T=8$-critical cells are closed
within the tested primes, $p\le 61$, by exact census plus the shear
step); \emph{rungs} (the frontier cells of every $T\ge 8$ are closed
only conditionally, by Theorem~\ref{thm:SC} under its ledger); and
\emph{rigid-zone finiteness} (open; the five-prime zone of
\S\ref{sec:rigid} is an empirical finding of the verified range).
Statements below use the matching noun.

\paragraph{Where the derivation lives.}
Section~\ref{sec:scale} derives the cascade model and proves the proved
pieces (Lemmas~\ref{lem:twofiber}, \ref{lem:pool}, \ref{lem:spacing},
\ref{lem:reflect}, \ref{lem:renorm}; Theorem~\ref{thm:chaindilate});
reports the measurements with their scopes and ground-truth records;
states the two open obligations (the strand census for H1g, the sharp
pair-line constants for H2m) in their reduced forms; and closes with
the epistemic status. Theorem~\ref{thm:SC} should be read as the
summary of that section, not as an independent result.
